%% ============================================================================
%%  back/appendix1.tex
%%  پیوست الف: اثبات قضیه‌ها و جزئیات محاسباتی
%%  طراحی شده برای قالب رسمی دانشگاه خوارزمی
%% ============================================================================
%%
%%  توضیح:
%%  این پیوست شامل اثبات‌های ریاضی تفصیلی، روابط تحلیلی، و جداول تکمیلی است
%%  که برای حفظ یکپارچگی و روانی متن اصلی در فصول، به بخش پیوست‌ها منتقل شده‌اند.
%% ============================================================================

\chapter{اثبات قضیه‌ها و جزئیات محاسباتی}
\label{app:proofs}

در این پیوست، اثبات‌های دقیق ریاضی، روابط تحلیلی همگرایی الگوریتم‌های پیشنهادی و همچنین جزئیات جدول فراسنج‌های محاسباتی تشریح می‌گردد.

\section{اثبات قضیه همگرایی کران‌دار}
\label{sec:app_convergence}

فرض کنید تابع هدف $f: \mathbb{R}^d \to \mathbb{R}$ تابعی محدب و مشتق‌پذیر با گرادیان لیپ‌شیتس به ثابت $L > 0$ باشد؛ یعنی به ازای هر $x, y \in \mathbb{R}^d$ داشته باشیم:
\begin{equation}
\label{eq:lipschitz}
\|\nabla f(x) - \nabla f(y)\| \le L \|x - y\|
\end{equation}

بر اساس لم نزول (\lr{Descent Lemma})، برای هر $x, y \in \mathbb{R}^d$ نابرابری زیر برقرار است:
\begin{equation}
\label{eq:descent_lemma}
f(y) \le f(x) + \langle \nabla f(x), y - x \rangle + \frac{L}{2} \|y - x\|^2
\end{equation}

اکنون قاعده به‌روزرسانی گرادیان با نرخ یادگیری $\eta$ را به صورت $x_{k+1} = x_k - \eta \nabla f(x_k)$ در نظر می‌گیریم. با جایگذاری $y = x_{k+1}$ در رابطه \eqref{eq:descent_lemma} به دست می‌آید:
\begin{align}
\label{eq:proof_step1}
f(x_{k+1}) &\le f(x_k) - \eta \|\nabla f(x_k)\|^2 + \frac{L \eta^2}{2} \|\nabla f(x_k)\|^2 \nonumber\\
&= f(x_k) - \eta \left(1 - \frac{L\eta}{2}\right) \|\nabla f(x_k)\|^2
\end{align}

چنانچه نرخ یادگیری در بازه $0 < \eta < \frac{2}{L}$ انتخاب شود، ضریب گرادیان اکیداً منفی بوده و دنباله مقادیر $\{f(x_k)\}$ یکنواخت نزولی و همگرا به نقطه کمینه سراسری خواهد بود.

\section{جدول مقادیر و فراسنج‌های تجربی}
\label{sec:app_params}

در جدول \ref{tab:app_hyperparams}، مقادیر اولیه و بازه جستجوی فراسنج‌های مدل بهینه‌سازی آورده شده است.

\begin{table}[H]
\centering
\caption{فراسنج‌های الگوریتم و تنظیمات شبیه‌سازی}
\label{tab:app_hyperparams}
\begin{tabular}{|l|c|c|p{6cm}|}
\hline
\textbf{فراسنج} & \textbf{نماد ریاضی} & \textbf{مقدار پیشنهادی} & \textbf{توضیحات و کاربرد} \\
\hline
نرخ یادگیری اولیه & $\eta_0$ & ۰/۰۰۱ & نرخ گام در گرادیان کاهشی \\
ضریب ممان & $\beta_1$ & ۰/۹ & ضریب میانگین متحرک نمایی مرتبه اول \\
ضریب مقیاس‌بندی & $\beta_2$ & ۰/۹۹۹ & ضریب میانگین مجذور گرادیان‌ها \\
ثابت پایداری عددی & $\epsilon$ & $10^{-8}$ & جلوگیری از تقسیم بر صفر در مخرج کسر \\
اندازه دسته داده & $B$ & ۶۴ & تعداد نمونه‌های مورد پردازش در هر تکرار \\
تعداد کل دورها & $E$ & ۱۰۰ & دورهای کامل آموزش روی مجموعه داده \\
\hline
\end{tabular}
\end{table}

\section{قطعه کد محاسبات کمکی}
\label{sec:app_code}

در ادامه، قطعه کد پیاده‌سازی تابع ارزیابی همگرایی در زبان پایتون نمایش داده شده است:

\begin{latin}
\begin{lstlisting}[language=Python, caption={پیاده‌سازی پایش همگرایی (\lr{Convergence Monitoring Implementation})}]
import numpy as np

def check_convergence(grad_norm, tol=1e-5, max_iter=1000):
    """
    Evaluates gradient descent stopping criterion.
    """
    for iteration, g in enumerate(grad_norm):
        if g < tol:
            return True, iteration
        if iteration >= max_iter:
            break
    return False, max_iter
\end{lstlisting}
\end{latin}
